\section{Discrete birth--death--catastrophe}
\label{m:app:dbdc}

This appendix develops the discrete-generation counterpart of the
birth--death--catastrophe process of \cref{m:def:bdc}: exact probabilities for
rupture and survival in discrete time, the weighted generating-function
recursion that carries the joint law, and the killed-generator formulation of
quasi-stationarity used by the continuous-time analysis.

\subsection*{Probability of catastrophe at step \texorpdfstring{$n$}{n}}

Let $\kappa$ be the probability that a given individual initiates catastrophe at
a step, and write $c = 1-\kappa$. A single individual fails to do so with
probability $c$, two with $c^2$, and a cohort of size $N$ with $c^N$.

Consider first birth and catastrophe alone, with no death, so $p=1$. Then the
census of the living population at any generation is known with certainty: from
a single progenitor the population simply doubles at each step prior to
catastrophe. Conditional on reaching step $n$, where $2^{n-1}$ individuals are
exposed, the process therefore comes to its abrupt halt at that step with
probability
\begin{equation}
\label{m:eq:bdchalt}
h_n = 1-c^{2^{\,n-1}},
\end{equation}
one minus the probability that none of the $2^{n-1}$ individuals present caused
it. This is a \emph{conditional} hazard: it assumes step $n$ is reached. The
unconditional probability that catastrophe first occurs at step $n$ is smaller,
because every earlier step must also have been passed.

\subsection*{Probability of process survival}

Still with no death, define $S_n$ as the probability that the process survives
to generation $n$. Survival to generation $0$ is certain, $S_0=1$. Survival to
generation $1$ requires only that the founder divided without catastrophe,
$S_1=c$. Persistence to step $2$ requires that the founder divided and that each
of its two offspring did too, so $S_2 = c\,c^2$. The pattern is clear: survival
to generation $n$ requires every prior step to have been passed, and therefore
\begin{equation}
\label{m:eq:bdcS}
S_n = \prod_{i=0}^{n-1}c^{2^i}
= c^{\sum_{i=0}^{n-1}2^i}
= c^{2^n-1} .
\end{equation}
The unconditional mass function of the first-catastrophe time follows at once:
\begin{equation}
\label{m:eq:bdcfirst}
\pr[\big]{K = n} = S_n\,h_n = c^{2^n-1}\lrb{1-c^{2^n}} .
\end{equation}
The factor $S_n$ is indispensable; $h_n$ alone is a conditional hazard, not an
exact-time probability, and conflating the two is the standard error of
mean-field hazard reasoning, which substitutes the current expected census into a
rate as though the census were deterministic.

\subsection*{Death included}

The difficulty comes once death is allowed, because the size of the living
population is then no longer known with certainty from one generation to the
next. Let $\widetilde Z_i$ denote the underlying binary Galton--Watson
population with the catastrophe mechanism switched off, and write
\begin{equation}
\label{m:eq:Hn}
H_n = \sum_{i=0}^{n-1}\widetilde Z_i
\end{equation}
for the cumulative exposure over the first $n$ steps. Conditioning on the whole
catastrophe-free genealogy gives the exact identity
\begin{equation}
\label{m:eq:noruptureexact}
\pr[\big]{\text{no catastrophe before }n}
= \ex{c^{H_n}} .
\end{equation}
The natural guess is that the probabilities shake out when the expected
population is used in place of the actual one:
\begin{equation}
\label{m:eq:bdcguess}
\ex{c^{H_n}} \overset{?}{=} c^{\sum_{i=0}^{n-1}\ex{\widetilde Z_i}} .
\end{equation}
This is a heuristic substitution and not an identity; the direction of the error
can be pinned down exactly. The map $N\mapsto c^N$ is convex, so Jensen's
inequality gives $\ex{c^{\widetilde Z_i}}\ge c^{\ex{\widetilde Z_i}}$: replacing
the random census by its mean systematically understates the chance of getting
through a step. With $m=2p$,
\begin{equation}
\label{m:eq:meanexposure}
\ex{H_n} = \sum_{i=0}^{n-1}m^i
= \begin{cases}
\dfrac{1-m^n}{1-m}, & m\neq1,\\[2mm]
n, & m=1,
\end{cases}
\end{equation}
and \cref{m:eq:bdcguess} is better read as the lower bound
$\pr[\big]{\text{no catastrophe before }n}\ge c^{\ex{H_n}}$ than as an equality.
Nor are the steps independent, since surviving one is evidence of a smaller
population, which raises the chance of surviving the next; both effects push the
same way. How tight a bound remains a question for numerical work.

Non-rupture should not be confused with population survival. If the requirement
is that the population be both unruptured and non-zero at generation $n$, the
exact probability is
\begin{equation}
\label{m:eq:jointsurvival}
\ex{\ind_{\{\widetilde Z_n>0\}}\,c^{H_n}},
\end{equation}
for which the Jensen bound does not directly provide an estimate. The
distinction removes an ambiguity in the elementary approximation and indicates
why the joint quasi-stationary problem is better posed through the
killed-chain generator developed below.

The process can nevertheless be described exactly by a weighted generating
function. Define
\begin{equation}
\label{m:eq:catweightedpgf}
F_n(s) = \ex{s^{\widetilde Z_n}\,c^{H_n}},
\qquad F_0(s) = s .
\end{equation}
Conditioning on the founder's reproduction --- death with probability $1-p$, or
division with probability $p$, after which two independent subtrees each
contribute a factor $F_n$ and the founder's own step contributes $c$ --- gives
the closed recursion
\begin{equation}
\label{m:eq:catweightedrec}
F_{n+1}(s) = c\,\phi\bigl(F_n(s)\bigr)
= c\lrb{1-p+pF_n(s)^2} .
\end{equation}
Its evaluations at special points have distinct meanings: $F_n(1)$ is the
probability of no catastrophe in the first $n$ updates, agreeing with
\cref{m:eq:noruptureexact}, while $F_n(1)-F_n(0)$ is the probability of being
both unruptured and non-extinct at generation $n$, the quantity of
\cref{m:eq:jointsurvival}. The recursion \cref{m:eq:catweightedrec} is exact,
and it is the natural object to iterate when the Jensen bound is too crude.

\subsection*{The killed-chain viewpoint}

The continuous-time analysis of birth--death--catastrophe processes proceeds
through the generator restricted to the living states, and it is useful to fix
the language here. Allow general state-dependent rates for the moment: from
state $i\ge1$, birth at rate $\lambda_i$, ordinary death at rate $\mu_i$, and
catastrophe at rate $\kappa_i$, with both absorbing outcomes lumped into a
single cemetery state. On bounded test functions the generator is
\begin{multline*}
(\mathcal Lf)(i)
= \lambda_i\{f(i+1)-f(i)\} + \mu_i\{f(i-1)-f(i)\}
+ \kappa_i\{f(0)-f(i)\},
\qquad i\ge1,
\end{multline*}
with $(\mathcal Lf)(0)=0$. Linear birth--death with per-capita catastrophe
corresponds to $\lambda_i=i\lambda$, $\mu_i=i\mu$, $\kappa_i=i\rho$, recovering
\cref{m:def:bdc} once the cemetery is split into extinction and rupture.

Restricting $\mathcal L$ to the positive states gives the \emph{killed
subgenerator} $K$, with
\begin{equation*}
K_{i,i+1}=\lambda_i,\qquad K_{i,i-1}=\mu_i\quad(i\ge2),
\qquad
K_{i,i}=-(\lambda_i+\mu_i+\kappa_i) .
\end{equation*}
The missing row mass is exactly the instantaneous absorption rate: $\kappa_i$
for $i\ge2$, and $\mu_1+\kappa_1$ for $i=1$, since from a single individual
ordinary death is also absorption. Writing
$T_0=\inf\{t\ge0:X_t=0\}$ and summing the positive-state forward equations gives
the survival-loss balance
\begin{equation}
\label{m:eq:survivalloss}
\dda{}{t}\pr{T_0>t}
= -\mu_1\pr{X_t=1} - \sum_{i\ge1}\kappa_i\pr{X_t=i},
\end{equation}
where ordinary death removes survival mass only from state $1$, while
catastrophe can remove it from every positive state.

The balance sharpens the connection with quasi-stationarity. If a distribution
$\nu=(\nu_i)_{i\ge1}$ on the positive states is quasi-stationary and the
relevant sums are finite, then for some $\theta>0$,
\begin{equation}
\label{m:eq:qsdeigen}
\nu K = -\theta\,\nu,
\qquad
\pr[\big]{_\nu}{T_0>t} = \ee^{-\theta t} .
\end{equation}
Indeed, quasi-stationarity writes the unnormalised positive-state law as
$a(t)\nu$; the semigroup property forces $a(s+t)=a(s)a(t)$, and
right-continuity gives $a(t)=\ee^{-\theta t}$; differentiating at zero yields
the left-eigenvector relation. Summing its coordinates and comparing with
\cref{m:eq:survivalloss} identifies the decay rate as
\begin{equation}
\label{m:eq:thetarate}
\theta = \mu_1\nu_1 + \sum_{i\ge1}\kappa_i\nu_i .
\end{equation}

One limiting case is worth isolating. If the catastrophe clock is constant,
$\kappa_i\equiv\kappa$ on the positive states, it may be taken independent of
the underlying birth--death process, and the killed semigroup factorises as
\begin{equation*}
P_t^{K} = \ee^{-\kappa t}\,P_t^{\mathrm{BD},K} .
\end{equation*}
The extra clock shifts the survival decay rate by $\kappa$ but leaves the law
conditional on survival unchanged. A state-dependent sequence such as
$\kappa_i=i\rho$, by contrast, weights paths according to their accumulated
population exposure and generally changes the quasi-stationary law itself.
Rates, means and variances, state probabilities and the associated
generating-function equations for \cref{m:def:bdc} are developed \fwd{bdc}{in
the later chapters on birth--death--catastrophe processes}, with the two-type
extension treated \fwd{multitype}{in the chapter on multi-type processes} and
the release of a ruptured compartment's contents into a shared medium
\fwd{rupture}{in the chapter on compartment rupture}.

\FloatBarrier
